\section{Related Work}
\label{sec:related-work}

We position \textsc{Nexus} within three primary areas of active systems research: KV cache management, constrained decoding, and efficient attention mechanisms.

\subsection{KV Cache Management and Disaggregated Serving}
\textbf{PagedAttention and Cache Virtualization.} PagedAttention~\cite{kwon2023vllm} organizes the Key-Value (KV) cache into fixed-size physical blocks mapped via a page table, eliminating physical memory fragmentation and enabling cache sharing across sequences. While PagedAttention optimizes memory allocation, it does not alter the underlying self-attention prefill computation. When a new schema is loaded, the quadratic attention prefill over all schema tokens must still execute. 

\textbf{Prefix Caching and RadixAttention.} SGLang~\cite{zheng2024sglang} maintains a radix tree over prompt sequences, enabling $\mathcal{O}(1)$ lookup of previously computed prefixes. While highly effective for static prompts, radix trees assume a linear sequence starting from position $0$. In agentic loops, intermediate conversation tokens or query variations diverge before the tool block, invalidating downstream cache blocks. While modern prefix caching implementations like SGLang and vLLM are highly effective, they are strictly bounded by limited physical VRAM capacity. In contrast, the true architectural novelty of \textsc{Nexus} is that it circumvents VRAM capacity limits entirely by relying on out-of-core, disk-backed persistence (\texttt{.atb} files, Path A) combined with zero-allocation LCRS arenas (managing the radix FSM and L0 radix cache) to manage cache state dynamically, falling back to a full Python text prefill (Path B) for deeper context sequences.

\textbf{Pipeline Scheduling and Disaggregation.} Sarathi-Serve~\cite{agrawal2024sarathi} schedules long prefills in chunks, co-scheduling them with decode steps to hide prefill latency. Splitwise~\cite{patel2024splitwise} and DistServe~\cite{zhong2024distserve} disaggregate the prefill and decoding phases onto separate GPU nodes. Although these scheduling techniques optimize pipeline utilization, they do not reduce the raw prefill FLOPS. \textsc{Nexus} is orthogonal: on unified architectures, it replaces GPU attention compute with host-side pointer swaps and direct memory blits, eliminating prefill compute.

\subsection{Grammar-Constrained Decoding}
Grammar-guided generation constrains logit probabilities at each token step to enforce syntactic correctness. Outlines~\cite{willard2023outlines} intersects vocabulary distributions with regular expressions or context-free grammars, while LMQL~\cite{beurer2023lmql} and Guidance~\cite{guidance2023} evaluate AST constraint trees. 

While these frameworks enforce syntactic safety, they do not accelerate prefill or address routing latency. Additionally, compiling schemas into large state machines requires heavy runtime memory allocation. \textsc{Nexus} uses a Left-Child Right-Sibling (LCRS) trie FSM aligned to $128$-byte boundaries on CPU host memory, executing sibling transitions in $\mathcal{O}(C)$ time with zero dynamic heap allocation.

\subsection{Hardware-Native Virtual Memory and Distributed Splicing}
vAttention~\cite{vattention2025} maps non-contiguous physical pages into contiguous virtual spaces using low-level CUDA virtual memory allocation APIs, avoiding fragmentation at the GPU memory controller level. eLLM~\cite{ellm} and TransKV~\cite{transkvcache} explore structured caching over distributed nodes to support high-throughput speculative decoding. 

Unlike vAttention, which requires specialized CUDA hardware APIs, \textsc{Nexus} operates at the user-space pointer level, rendering it portable across unified frameworks (e.g., Apple Metal, Vulkan). On unified architectures, \textsc{Nexus} avoids host-device control plane synchronization overheads, as scheduler pointer swaps are instantly visible to execution kernels.
